% !TeX root = ../main.tex

\chapter{结论与展望}

本文围绕多中心海洋浮游生物监测场景下“原始图像难以跨中心共享，而监测系统又需要具备 OOD 拒识能力”的问题，研究了联邦后处理 OOD 检测方法。针对这一问题，本文提出了 FedViM，并在此基础上进一步构造了其自适应选维扩展 ACT-FedViM。其中，FedViM 通过联邦聚合一阶与二阶充分统计量，实现了 ViM 所需全局特征均值与协方差的重构；ACT-FedViM 则在此基础上利用 ACT 自动选择主子空间维度 $k$，从而形成一套统计驱动的联邦 ViM 后处理方法。

在基于 DYB-PlanktonNet 构建的 OOD 划分上，本文对 5 个 CNN 模型进行了实验评估。结果表明，联邦化 ViM 路线在该任务上是可行的：五模型平均 ID 分类准确率达到 96.48\%，说明联邦分类模型具备较好的已知类识别能力；与 MSP 和 Energy 两类输出空间基线相比，FedViM 与 ACT-FedViM 在 Near-OOD 和 Far-OOD 检测上整体表现更优，说明利用特征空间结构信息的后处理方法更适合当前细粒度浮游生物 OOD 场景。

综合全文的理论分析与实验结果，本文得到以下几点主要结论：
\begin{enumerate}
  \item ViM 所需的全局特征统计量可以通过联邦充分统计量聚合完成重构，从而形成可实现的联邦后处理 OOD 检测流程。
  \item ACT 可以作为联邦 ViM 的自适应选维模块接入已经训练好的分类模型的后处理评估，为主子空间维度选择提供数据驱动的统计依据。
  \item ACT-FedViM 的主要价值并不完全体现在 AUROC 提升上，而更多体现在统计驱动的选维方式、压缩后处理主子空间规模以及提升部署友好性等方面。
\end{enumerate}

总体而言，FedViM 给出了联邦场景下实现 ViM 后处理 OOD 检测的可行方案，ACT-FedViM 则在此基础上进一步提高了后处理模块的紧凑性与对不同模型的适应性。

尽管本文在联邦后处理 OOD 检测方面进行了初步探索，但当前工作仍存在一定局限。
\begin{enumerate}
  \item 本文对 ViM 的改进主要体现在主子空间维度 $k$ 的选择上，没有改变主子空间与残差空间的构造方式。因此，当 OOD 判别信息没有落在样本协方差矩阵的残差空间时，自适应选维所能带来的性能改善仍然是有限的。未来可在保持联邦统计量可聚合的前提下，进一步研究 OOD 判别信息的结构，设计出能够涵盖更多 OOD 判别信息的选维机制。
  \item 本文的方法虽然避免了原始图像和逐样本特征的跨中心传输，但在传输统计量时并没有引入隐私机制来给出数学上的隐私保证。因此，后续研究可在统计量聚合流程中引入安全聚合、差分隐私扰动或其他隐私保护机制，进一步提升多中心海洋浮游生物监测系统的隐私保护能力。
\end{enumerate}
